\title{Group Sequence Policy Optimization}

\begin{abstract}
We present Group Sequence Policy Optimization (GSPO), a robust, effective, and high-performing reinforcement learning algorithm specifically developed for training large language models. In contrast to prior methods that rely on token-level importance ratios, GSPO constructs the importance ratio at the sequence level, applying sequence-based clipping, reward assignment, and optimization processes. Our experiments reveal that GSPO surpasses the GRPO algorithm in terms of training efficiency and overall performance, delivers enhanced stability for Mixture-of-Experts (MoE) reinforcement learning training, and offers potential to streamline RL infrastructure development. The strengths of GSPO have played a key role in the significant advancements realized in the most recent Qwen3 model releases.
\end{abstract}

\section{Introduction}
\label{sec:intro}

Reinforcement learning (RL) has become a central framework for advancing the scalability of language models \citep{o1,dpsk-r1,qwq32b,qwen3}. By leveraging RL at scale, language models acquire the proficiency to address complex tasks—including advanced mathematics and programming—by engaging in more extended and intricate reasoning sequences.

Achieving scalable reinforcement learning (RL) through increased computational resources fundamentally depends on ensuring stable and resilient training behavior. Nevertheless, leading RL approaches such as GRPO \citep{grpo} demonstrate significant stability challenges when applied to the training of large-scale language models, frequently leading to disastrous and unrecoverable collapse of the models \citep{qwen3, minimax-m1}. These persistent instabilities impede attempts to further advance language model performance via extended RL-based optimization.

In this work, we demonstrate that GRPO's instability originates from the incorrect use and breakdown of importance sampling weights within its algorithmic structure. This leads to substantial training variance, which grows as response sequences become longer and is exacerbated by the implementation of the clipping method, eventually causing the model to fail.

In response to these fundamental challenges, we introduce \textbf{Group Sequence Policy Optimization (GSPO)}, a novel reinforcement learning framework tailored for large language model training. GSPO’s principal contribution is its rigorous formulation of the importance ratio derived from sequence likelihood, adhering to the essential tenets of importance sampling as described by \citep{click}. Moreover, GSPO evaluates normalized rewards as the advantages obtained across several generated responses for each query, thereby enhancing the coherence between sequence-wise reward assignment and the optimization process.

Our experimental results reveal that GSPO outperforms GRPO considerably in terms of training stability, effectiveness, and overall results. Notably, GSPO naturally addresses the stability problems commonly encountered during RL training of large Mixture-of-Experts (MoE) architectures, removing reliance on elaborate stabilization mechanisms and indicating possibilities for streamlined RL systems. The strengths of GSPO have played a key role in the remarkable advancements seen in recent Qwen3 models. We anticipate that GSPO will serve as a reliable and extensible foundation for the ongoing progress of large-scale RL-based language model training.

\section{Preliminaries}

\paragraph{Notation}
Throughout this work, we represent an autoregressive language model with parameters $\theta$ as a policy, denoted by $\pi_\theta$.
The symbol $x$ signifies a query, while $\mathcal{D}$ stands for the set of all queries. For any response $y$ corresponding to a query $x$, its probability under the policy $\pi_\theta$ is given by $\pi_\theta (y | x)=\prod_{t=1}^{|y|} \pi_\theta (y_t | x, y_{<t} )$, with $|y|$ indicating the length of $y$ in tokens. Each query-response pair $(x, y)$ receives an evaluation from a verifier $r$, which assigns a reward $r(x, y) \in [0, 1]$.

\paragraph{Proximal Policy Optimization (PPO)} Proximal Policy Optimization (PPO) \citep{ppo} leverages data collected from the prior policy $\pi_{\theta_\text{old}}$, ensuring that updates to the policy remain close to this baseline via a clipping strategy. More precisely, PPO maximizes the following surrogate objective to optimize the policy (for clarity and because it is not central to this work, we exclude the KL regularization term in what follows):

\begin{align}
\mathcal{J}_\text{PPO}(\theta) = \mathbb{E}_{ x \sim \mathcal{D},\, y \sim \pi_{\theta_\text{old}}( \cdot | x) }
\left[ \frac{1}{|y|} \sum_{t=1}^{|y|} 
\min \left( w_{t}(\theta) \widehat{A}_{t},  \, \mathrm{clip} \left( w_{t}(\theta), 1 - {\varepsilon}, 1 + {\varepsilon}\right) \widehat{A}_{t} \right)
\right],
\end{align}

Here, the token $y_t$'s importance ratio is specified by $
w_{t}(\theta) = \frac{ \pi_{\theta} (y_{t} | x, y_{<t}) }{ \pi_{\theta_\text{old}} (y_{t} | x,y_{<t})} $. The advantage $\widehat{A}_{t}$ for $y_t$ is obtained through a separate value function, and $\varepsilon$ denotes the threshold for clipping the importance ratios.

A fundamental difficulty encountered with PPO in real-world applications is its substantial dependence on the value model. Typically, the value model possesses a complexity and scale comparable to that of the policy model, which significantly escalates both memory usage and computational costs. Additionally, the performance of the algorithm is closely tied to the accuracy of the value estimation. Developing a trustworthy value model is intrinsically difficult, but extending its reliability to accommodate longer outputs and increasingly sophisticated tasks is an even more daunting problem.

\paragraph{Group Relative Policy Optimization (GRPO)} GRPO \citep{grpo} eliminates reliance on a value model by directly evaluating the relative advantages among a collection of responses associated with an identical query. In particular, GRPO seeks to maximize the following objective:

\begin{align}
\label{equ:grpo}
\mathcal{J}_\text{GRPO}(\theta) = \mathbb{E}_{ x \sim \mathcal{D},\, \{y_i\}_{i=1}^G \sim \pi_{\theta_\text{old}}( \cdot | x) }
\left[ \frac{1}{G} \sum_{i=1}^{G} \frac{1}{|y_i|} \sum_{t=1}^{|y_i|} 
\min \left( w_{i,t}(\theta) \widehat{A}_{i,t},  \, \mathrm{clip} \left( w_{i,t}(\theta), 1 - {\varepsilon}, 1 + {\varepsilon}\right) \widehat{A}_{i,t} \right)
\right],
\end{align}

Here, $G$ denotes the total number of responses produced for each query $x$ (referred to as the group size). The importance ratio $w_{i,t}(\theta)$ and the estimated advantage $\widehat{A}_{i,t}$ corresponding to token $y_{i,t}$ are defined as follows:

\begin{align}
    w_{i,t}(\theta)=\frac{ \pi_{\theta} (y_{i,t} | x, y_{i,<t}) }{ \pi_{\theta_\text{old}} (y_{i,t} | x,y_{i,<t})},\quad\ 
    \widehat{A}_{i,t} = \widehat{A}_{i} = \frac{r(x, y_i) - \mathrm{mean} \left( \{ r(x, y_i) \}_{i=1}^G \right) }{ \mathrm{std} \left( \{ r(x, y_i) \}_{i=1}^G \right) },
\end{align}

respectively, with every token in $y_i$ assigned the identical advantage value $\widehat{A}_{i}$.

\section{Motivation}
\label{sec:motivation}

As model sizes increase and architectures incorporate greater sparsity (such as in Mixture-of-Experts models), along with longer generated responses, achieving efficient use of computational resources in reinforcement learning requires scaling up rollout batch sizes. To enhance sample efficiency, practitioners typically subdivide these extensive rollout batches into several mini-batches for the purposes of gradient computation. However, this batching approach gives rise to an off-policy learning scenario: the sampled outputs $y$ are generated using a previous policy $\pi_{\theta_\text{old}}$, not the updated policy $\pi_\theta$ that is currently subject to optimization. This off-policy nature justifies the presence of clipping strategies in methods like PPO and GRPO, which act to restrict the influence of data points that significantly diverge from the current policy distribution during gradient calculation.

Although techniques such as clipping are used to address off-policy discrepancies, we argue that GRPO suffers from a deeper flaw: \textit{its objective is fundamentally ill-posed}. This issue is especially problematic in the context of training large-scale models for tasks involving lengthy responses, where it can result in severe model failure. The root cause of this ill-posedness in the GRPO objective lies in an inappropriate use of importance sampling weights. Traditionally, importance sampling is intended to estimate the expectation of a function $f$ under a target distribution $\pi_\text{tar}$ by adjusting the weights of samples collected from a behavior distribution $\pi_\text{beh}$ as follows:

\begin{align}
\mathbb{E}_{z \sim \pi_\text{tar}} \left[ f(z) \right] = \mathbb{E}_{z \sim \pi_\text{beh}} \left[ \frac{ \pi_\text{tar} (z) }{ \pi_\text{beh} (z)} \, f(z) \right].
\end{align}

Importantly, effective correction for the distributional discrepancy using the importance weight $\frac{ \pi_\text{tar} (z) }{ \pi_\text{beh} (z)}$ depends on averaging across a large number of samples ($N \gg 1$) drawn from the behavior distribution $\pi_\text{beh}$.

On the other hand, GRPO implements the importance weight $\frac{ \pi_{\theta} (y_{i,t} | x, y_{i,<t}) }{ \pi_{\theta_\text{old}} (y_{i,t} | x,y_{i,<t})}$ at every token index $t$. This approach, relying on a solitary sample $y_{i,t}$ drawn from each corresponding next-token distribution $\pi_{\theta_\text{old}}( \cdot | x, y_{i, <t})$, is insufficient for proper distribution correction. Rather than achieving its intended purpose, this method infuses significant variance into the gradient updates, particularly across longer sequences—a phenomenon that is further intensified by the presence of clipping. Our experimental results indicate that such training instability can induce irreversible model collapse. After such collapse, attempts to restore effective training—whether by reverting to earlier checkpoints, carefully adjusting hyperparameters (such as clipping thresholds), increasing output sequence length, or altering the reinforcement learning queries—prove futile.

The aforementioned finding highlights a key flaw inherent in GRPO’s architecture. Specifically, the ineffectiveness of token-level importance weighting underscores a central tenet: \textbf{the granularity of the optimization objective must correspond to that of the reward signal}. Given that the reward is attributed to whole sequences, implementing off-policy corrections at the individual token level may be inappropriate. This realization leads us to abandon token-level objectives in favor of approaches that incorporate importance weighting and carry out optimization at the \textit{sequence level} instead.

\section{Algorithm}

\subsection{GSPO: Group Sequence Policy Optimization}

Although GRPO faces challenges with token-level importance weighting, specifically $\frac{ \pi_{\theta} (y_{i,t} | x, y_{i,<t}) }{ \pi_{\theta_\text{old}} (y_{i,t} | x,y_{i,<t})}$, it is noteworthy that, in language generation tasks, the \textit{sequence-level} importance weight $\frac{ \pi_{\theta} (y | x) }{ \pi_{\theta_\text{old}} (y | x)}$ possesses a well-defined theoretical significance. This ratio captures the extent to which the response $y$, drawn from $\pi_{\theta_\text{old}} (\cdot | x)$, diverges from the newer model $\pi_{\theta} (\cdot | x)$. Consequently, this sequence-level metric naturally corresponds to the sequence-level reward and functions as an informative measure for the clipping operation.

Drawing from this simple insight, we introduce the \textbf{Group Sequence Policy Optimization (GSPO)} algorithm. GSPO utilizes the sequence-level optimization goal outlined below:

\begin{align}
\mathcal{J}_\text{GSPO} (\theta) =
\mathbb{E}_{ x \sim \mathcal{D},\, \{y_i\}_{i=1}^G \sim \pi_{\theta_\text{old}}( \cdot | x) }
\left[ 
\frac{1}{G} \sum_{i=1}^{G}
\min \left( s_{i}(\theta)  \widehat{A}_{i},  \, \mathrm{clip} \left( s_{i}(\theta), 1 - {\varepsilon}, 1 + {\varepsilon} \right) \widehat{A}_{i} \right) 
\right],
\label{equ:gspo}
\end{align}

in which we utilize the group-wise approach for estimating advantage:

\begin{align}
\widehat{A}_{i} = \frac{r(x, y_i) - \mathrm{mean} \left( \{ r(x, y_i) \}_{i=1}^G \right) }{ \mathrm{std} \left( \{ r(x, y_i) \}_{i=1}^G \right) },
\end{align}

and establish the importance ratio $s_{i}(\theta)$ using the probability of the sequence \citep{click}:

\begin{align}
s_{i}(\theta) = \left( \frac{ \pi_{\theta} (y_i | x) }{ \pi_{\theta_\text{old}} (y_i | x)} \right)^{\frac{1}{|y_i|}}
=
\exp \left( \frac{1}{|y_i|} \sum_{t=1}^{|y_i|} \log \frac{ \pi_{\theta} (y_{i,t} | x, y_{i,<t}) }{ \pi_{\theta_\text{old}} (y_{i,t} | x,y_{i,<t})} \right).
\end{align}

Consequently, GSPO conducts clipping at the response level rather than at the token level, aiming to filter out excessively ``off-policy'' samples from the gradient estimation process. This approach aligns the optimization procedure with the sequence-level reward structure. Additionally, length normalization is utilized in $s_{i}(\theta)$, serving to lower the variance and ensure that $s_{i}(\theta)$ remains within a consistent numerical scale. Without this normalization, significant alterations in the likelihoods of just a few tokens could cause large fluctuations in the overall importance ratio at the sequence level, necessitating different clipping thresholds for responses of various lengths. It is also worth mentioning that the typical clipping ranges used in GSPO and in earlier algorithms (such as GRPO) often differ by several orders of magnitude, attributable to the differences in how the importance ratios are formulated in these methods.

\subsection{Gradient Analysis}
\label{subsec:gradient}

The gradient of the GSPO objective can be computed in the following manner (for simplicity, we exclude clipping):

\begin{align}
\nabla_{\theta} \mathcal{J}_\text{GSPO} (\theta)
=&\ 
\nabla_{\theta} \mathbb{E}_{ x \sim \mathcal{D},\, \{y_i\}_{i=1}^G \sim \pi_{\theta_\text{old}}( \cdot | x) }
\left[ 
\frac{1}{G} \sum_{i=1}^{G} s_{i}(\theta) \widehat{A}_{i}
\right] \\
= &\ 
\mathbb{E}_{ x \sim \mathcal{D},\, \{y_i\}_{i=1}^G \sim \pi_{\theta_\text{old}}( \cdot | x) }
\left[ 
\frac{1}{G} \sum_{i=1}^{G}
s_{i}(\theta) \widehat{A}_{i}
\cdot \nabla_{\theta} \log s_{i}(\theta)
\right] \\
=&\ 
\mathbb{E}_{ x \sim \mathcal{D},\, \{y_i\}_{i=1}^G \sim \pi_{\theta_\text{old}}( \cdot | x) }
\left[ 
\frac{1}{G} \sum_{i=1}^{G} 
\left( \frac{ \pi_{\theta} (y_i | x) }{ \pi_{\theta_\text{old}} (y_i | x)} \right)^{\frac{1}{|y_i|}} \widehat{A}_{i} 
\cdot \frac{1}{|y_i|} \sum_{t=1}^{|y_i|} \nabla_{\theta} \log \pi_{\theta} (y_{i,t} | x, y_{i,<t}) 
\right].
\label{equ:gspo_grad}
\end{align}

To facilitate comparison, the gradient expression for the GRPO objective can be written as follows (where we set $\widehat{A}_{i,t} = \widehat{A}_{i}$):

\begin{align}
\nabla_{\theta} \mathcal{J}_\text{GRPO} (\theta) 
=&\ 
\nabla_{\theta} \mathbb{E}_{ x \sim \mathcal{D},\, \{y_i\}_{i=1}^G \sim \pi_{\theta_\text{old}}( \cdot | x) }
\left[ \frac{1}{G} \sum_{i=1}^{G} \frac{1}{|y_i|} \sum_{t=1}^{|y_i|} 
w_{i,t}(\theta) \widehat{A}_{i,t}
\right] \\
=&\
\mathbb{E}_{ x \sim \mathcal{D},\, \{y_i\}_{i=1}^G \sim \pi_{\theta_\text{old}}( \cdot | x) }
\left[ \frac{1}{G} \sum_{i=1}^{G} \widehat{A}_{i} 
\cdot \frac{1}{|y_i|} \sum_{t=1}^{|y_i|} 
\frac{ \pi_{\theta} (y_{i,t} | x, y_{i,<t}) }{ \pi_{\theta_\text{old}} (y_{i,t} | x,y_{i,<t})} 
\nabla_{\theta} \log \pi_{\theta} (y_{i,t} | x, y_{i,<t})  
\right].
\end{align}

The key difference between GSPO and GRPO pertains to \textit{the manner in which gradients of token log likelihoods are assigned weights}. In GRPO, each token receives a weight based on its "importance weight" $\frac{ \pi_{\theta} (y_{i,t} | x, y_{i,<t}) }{ \pi_{\theta_\text{old}} (y_{i,t} | x,y_{i,<t})}$. These non-uniform weights, which may fall within the interval $(0, 1+\varepsilon]$ when $\widehat{A}_{i} >0$, or within $[1-\varepsilon, +\infty)$ when $\widehat{A}_{i} < 0$, are significant and can aggregate over time, potentially resulting in unstable and unpredictable outcomes during training. Conversely, GSPO assigns identical weights to all tokens within a response, thus removing the source of instability introduced by the variable weighting scheme in GRPO.

\subsection{GSPO-token: A Token-level Objective Variant}

In contexts such as multi-turn RL, a more detailed advantage modification than what is possible at the sequence level might be preferable. To address this, we propose a token-level adaptation of GSPO, referred to as \textbf{GSPO-token}, which facilitates advantage tailoring at the granularity of individual tokens:

\begin{align}
\mathcal{J}_\text{GSPO-token}(\theta)
=
\mathbb{E}_{ x \sim \mathcal{D},\, \{y_i\}_{i=1}^G \sim \pi_{\theta_\text{old}}( \cdot | x) }
\left[ 
\frac{1}{G} \sum_{i=1}^{G} \frac{1}{|y_i|} \sum_{t=1}^{|y_i|} 
\min \left( s_{i,t}(\theta) \widehat{A}_{i,t},  \, \mathrm{clip} \left( s_{i,t}(\theta), 1 - {\varepsilon}, 1 + {\varepsilon} \right) \widehat{A}_{i,t} \right) 
\right],
\label{equ:gspo-token}
\end{align}

in which

\begin{align}
s_{i,t}(\theta) 
= \mathrm{sg} \left[ s_{i}(\theta) \right]  \cdot \frac{ \pi_{\theta} (y_{i,t} | x, y_{i,<t}) }{ \mathrm{sg} \left[ \pi_{\theta} (y_{i,t} | x, y_{i,<t}) \right] },
\end{align}

Here, $\mathrm{sg}[\cdot]$ refers to extracting just the numerical value while blocking any gradient flow, which is equivalent to the \texttt{detach} function in PyTorch. The gradient for GSPO-token can be calculated as follows:

\begin{align}
\nabla_{\theta} \mathcal{J}_\text{GSPO-token}(\theta)
=&\ 
\nabla_{\theta} \mathbb{E}_{ x \sim \mathcal{D},\, \{y_i\}_{i=1}^G \sim \pi_{\theta_\text{old}}( \cdot | x) }
\left[ 
\frac{1}{G} \sum_{i=1}^{G}
\frac{1}{|y_i|} \sum_{t=1}^{|y_i|} 
s_{i,t}(\theta) \widehat{A}_{i,t}
\right] \\
=&\ 
\mathbb{E}_{ x \sim \mathcal{D},\, \{y_i\}_{i=1}^G \sim \pi_{\theta_\text{old}}( \cdot | x) }
\left[ 
\frac{1}{G} \sum_{i=1}^{G} s_i (\theta) \cdot \frac{1}{|y_i|} \sum_{t=1}^{|y_i|} 
\widehat{A}_{i,t} \frac{ \nabla_{\theta} \pi_{\theta} (y_{i,t} | x, y_{i,<t}) }{ \pi_{\theta} (y_{i,t} | x, y_{i,<t}) }
\right] \\
=&\ 
\mathbb{E}_{ x \sim \mathcal{D},\, \{y_i\}_{i=1}^G \sim \pi_{\theta_\text{old}}( \cdot | x) }
\left[ 
\frac{1}{G} \sum_{i=1}^{G}  \left( \frac{ \pi_{\theta} (y_i | x) }{ \pi_{\theta_\text{old}} (y_i | x)} \right)^{\frac{1}{|y_i|}}
\cdot \frac{1}{|y_i|} \sum_{t=1}^{|y_i|}
\widehat{A}_{i,t} \nabla_{\theta} \log \pi_{\theta} (y_{i,t} | x, y_{i,<t})  
\right].
\label{equ:gspo-token_grad}
\end{align}

Observe that the factor $\frac{ \pi_{\theta} (y_{i,t} | x, y_{i,<t}) }{ \mathrm{sg} \left[ \pi_{\theta} (y_{i,t} | x, y_{i,<t}) \right] }$ evaluates to 1, rendering $s_{i,t}(\theta)$ and $s_{i}(\theta)$ equivalent in terms of their numerical values. By juxtaposing Equation~\eqref{equ:gspo} with~\eqref{equ:gspo-token}, as well as Equation~\eqref{equ:gspo_grad} with~\eqref{equ:gspo-token_grad}, one can see that GSPO-token and GSPO are numerically indistinguishable with respect to the optimization objective, the clipping mechanism, and the formal gradient—provided that all tokens in the sequence $y_i$ share a uniform advantage, that is, when $\widehat{A}_{i,t} = \widehat{A}_{i}$ for all $t$. Nevertheless, GSPO-token offers greater adaptability, as it allows the assignment of distinct advantage values to individual tokens.

\section{Experiments and Discussion}

\subsection{Empirical Results}

Our experiments utilize a cold-start model fine-tuned from Qwen3-30B-A3B-Base. We present both the reward trajectories during training and the resulting model performance across several benchmarks: AIME'24 (measured as the average Pass@1 across 32 samples), LiveCodeBench (from 202410 to 202502, evaluated as the average Pass@1 over 8 samples), and CodeForces (using Elo Rating). For reinforcement learning, each rollout batch is divided into four mini-batches to perform gradient updates. In the context of GSPO, the clipping intervals in Equation~\eqref{equ:gspo} are configured to 3e-4 (left) and 4e-4 (right). For comparison, we adopt GRPO as a baseline method, with its clipping parameters in Equation~\eqref{equ:grpo} set to 0.2 and 0.27, after extensive tuning to maintain fair evaluation conditions. It should be noted that GRPO requires the Routing Replay training strategy to achieve stable convergence for MoE RL, as discussed further in \S~\ref{subsec:routing_replay}, whereas \textit{GSPO eliminates the dependence on this strategy}.

As illustrated in Figure~\ref{fig:results}, training with GSPO exhibits consistent stability. Notably, \textbf{GSPO facilitates ongoing enhancement of model performance by allocating more training compute, periodically refreshing the query set, and lengthening generation output}. In addition, GSPO outperforms GRPO in terms of training efficiency, attaining higher accuracy and superior benchmark results while utilizing equivalent training compute and query resources. Finally, deploying GSPO in RL training for the recent Qwen3 models has yielded significant success, providing robust evidence for GSPO’s capability to harness the scaling potential of RL in large language model development.

\subsection{Curious Observation on Clipping Fractions}

One of the primary differences between GSPO and GRPO lies in their respective clipping strategies: GSPO applies clipping at the sequence level, truncating entire responses, whereas GRPO performs clipping at the level of individual tokens. As illustrated in Figure~\ref{fig:clipping}, this results in a stark discrepancy, with GSPO exhibiting a two-order-of-magnitude lower proportion of unclipped tokens relative to GRPO; adjusting the clipping thresholds does not alter this scale of difference. Interestingly, despite GSPO discarding a greater number of tokens—and therefore having fewer tokens available for training or for calculating gradients—it nonetheless attains greater training efficiency than GRPO. This unexpected outcome, wherein extensive token clipping is accompanied by improved training efficiency, suggests that GRPO’s approach of estimating gradients at the token level introduces substantial noise and diminishes sample efficiency. On the other hand, the sequence-level methodology used by GSPO yields a more robust and informative learning signal.

\subsection{Benefit of GSPO for MoE Training}
\label{subsec:routing_replay}

\paragraph{Background}
Relative to reinforcement learning (RL) training of conventional dense models, Mixture-of-Experts (MoE) architectures, due to their sparse activation patterns, face distinctive stability issues. Notably, when employing the GRPO algorithm, MoE models are susceptible to \textit{expert-activation volatility}, which can significantly impede the convergence of RL training. Specifically, after performing one or more updates of the model parameters, the composition of experts engaged for identical responses can shift markedly. For instance, in the case of the 48-layer Qwen3-30B-A3B-Base model, each RL gradient update leads to approximately a 10\% change in the set of experts activated for the same rollout sample when comparing the current policy $\pi_{\theta}$ to the previous policy $\pi_{\theta_\text{old}}$.
This instability, which tends to intensify in larger and deeper MoE configurations, results in large fluctuations in the token-level importance ratios $w_{i,t}(\theta) = \frac{ \pi_{\theta} (y_{i,t} | x, y_{i,<t}) }{ \pi_{\theta_\text{old}} (y_{i,t} | x,y_{i,<t})}$. Such volatility further undermines the reliability of these ratios, as elaborated in \S~\ref{sec:motivation} and~\ref{subsec:gradient}, and ultimately disrupts the normal progression towards convergence in RL training.

\paragraph{Our Previous Approach} In our earlier work, we addressed this issue using the \textbf{Routing Replay} training method. This involved recording the experts selected by $\pi_{\theta_\text{old}}$ and subsequently reusing, or ``replaying,'' these expert selections in $\pi_{\theta}$ when calculating the importance ratios $w_{i,t}(\theta) = \frac{ \pi_{\theta} (y_{i,t} | x, y_{i,<t}) }{ \pi_{\theta_\text{old}} (y_{i,t} | x,y_{i,<t})}$. Through this procedure, both $\pi_{\theta} (y_{i,t} | x, y_{i,<t})$ and $\pi_{\theta_\text{old}} (y_{i,t} | x,y_{i,<t})$ rely on the identical configuration of experts for each token $y_{i,t}$, allowing the token-level importance ratios to remain stable. This mechanism also guarantees that optimization is consistently applied to the same activated network during the process of gradient updates. As depicted in Figure~\ref{fig:routing_replay}, Routing Replay is a crucial mechanism for achieving stable convergence in the GRPO training of MoE architectures.

\paragraph{Benefit of GSPO} While Routing Replay allows MoE models trained with GRPO to reach convergence, it does so by reapplying prior routing patterns, which introduces both extra memory consumption and increased communication costs, while simultaneously restricting the effective capacity of the MoE architecture. On the other hand, as demonstrated in Figure~\ref{fig:results}, GSPO avoids reliance on Routing Replay entirely. It can straightforwardly compute the importance ratios $s_i(\theta)$, ensure proper convergence, and maintain stable optimization behavior. The fundamental observation behind this is that GSPO evaluates only the sequence likelihood (specifically, $\pi_{\theta} (y_i | x)$), disregarding the more granular likelihoods for individual tokens (such as $\pi_{\theta} (y_{i,t} | x, y_{i,<t})$). Since the language modeling competence of the MoE is always preserved, these sequence-level likelihoods remain comparatively stable. Overall, GSPO effectively addresses the instability associated with expert activation in MoE models, making complex mechanisms like Routing Replay unnecessary. As a result, it streamlines and stabilizes the training process and enables the model to fully utilize its representational capacity without imposed limitations.

\subsection{Benefit of GSPO for RL Infrastructure}

Due to differences in numerical precision between training engines (such as Megatron) and inference engines (like SGLang and vLLM), it is common in practical settings to recalculate the likelihoods of sampled outputs under the previous policy $\pi_{\theta_\text{old}}$ using the training engine. Nevertheless, GSPO performs optimization based on sequence-level likelihoods rather than on a token-level basis. Intuitively, sequence-level likelihoods are significantly more robust to precision differences. Consequently, GSPO allows the use of likelihoods generated directly from the inference engine for optimization, eliminating the requirement to recompute them with the training engine. This capability is particularly advantageous in contexts such as partial rollouts, multi-turn reinforcement learning, and frameworks where training and inference are decoupled.

\section{Conclusion}

We introduce Group Sequence Policy Optimization (GSPO), a novel reinforcement learning approach designed for training large language models. Building upon the foundation of importance sampling, GSPO establishes importance ratios grounded in sequence likelihood, and incorporates mechanisms for clipping, reward assignment, and optimization at the sequence level. Empirical results indicate that GSPO offers significantly enhanced training stability, efficiency, and overall performance in comparison to GRPO, and proves especially effective for large-scale reinforcement learning applied to MoE models. This algorithm underpins the substantial advancements observed in the most recent Qwen3 models. By adopting GSPO as a scalable methodological framework, we aim to further advance the scaling of reinforcement learning and anticipate consequential breakthroughs in the development of intelligence.

\end{document}